\chapter{Conclusion and Future Work}
\label{ch:conclusion}

This chapter summarizes the contributions of this thesis and outlines promising directions for future research in classical guitar playability analysis.

\section{Summary of Contributions}
This thesis has successfully adapted and evaluated the interpretable RubricNet architecture for the task of classical guitar difficulty estimation from symbolic notation. Our primary contributions are:
\begin{enumerate}
    \item \textbf{Explainable Guitar Architecture}: We deployed the first multi-dimensional, additive neural network model for classical guitar difficulty estimation, providing teachers and students with transparent, feature-level "rubric" score sheets for any graded piece.
    \item \textbf{Rhythm-Alignment Pipeline}: We implemented a Needleman-Wunsch sequence-matching pipeline to transplant MIDI temporal structures onto PDF-derived pitch vectors, resolving missing duration data for 89.2\% of our 716-piece dataset.
    \item \textbf{Feature Engineering Evolution}: We designed, audited, and refined three generations of symbolic descriptors. By adding fretboard coverage metrics (V2) and context-aware velocity/sliding-window bottleneck features (V3), we resolved key pedagogical limitations of piano-focused models.
    \item \textbf{Empirical Validation}: We evaluated our framework against the GuitarBurst graded database. RubricNet V3 achieved an exact 8-class accuracy of 31.0\% and a 3-class coarse accuracy of 67.6\%, outperforming a black-box Random Forest baseline across key ordinal metrics (balanced accuracy, MAE, MSE, and Kendall's $\tau$) while maintaining complete interpretability.
\end{enumerate}

\section{Future Work}
Several promising avenues remain to extend this research:

\subsection{Inferred Barre Chord Detection}
While the current dataset relies partly on transcription metadata to identify barre chords, many symbolic files omit these labels. A key future direction is implementing a geometric barre inference algorithm directly in the feature extractor. A chord event can be flagged as a barre if it contains $\ge 3$ fretted notes and at least 3 of those notes are played on the same fret across adjacent strings. The lowest string playing that fret acts as the barre root. This logic would allow the extraction of an \texttt{inferred\_barre\_ratio} descriptor from any clean score without transcription tags.

\subsection{Generative Rhythm Feature Imputation}
Currently, the 10.8\% of pieces lacking matched MIDI files receive median-imputed values for rhythm-dependent descriptors. A more advanced approach would train a generative model (such as a variational autoencoder or multi-task regressor) on the rhythm-complete subset to predict and impute missing temporal features directly from pitch-only sequences.

\subsection{Multi-Voice Separation Solvers}
To address the single-voice limitation of the alignment pipeline, future work should integrate multi-voice separation algorithms (such as the voice-separation tools in \texttt{music21}). Splitting the score into distinct treble, middle, and bass voices before alignment would enable more precise measurements of contrapuntal complexity.

\subsection{Fingering Solver Integration}
Unlike piano, where the staff layout strongly implies fingering, a single pitch on the guitar can be played at multiple fretboard positions. Integrating a physical guitar-fingering solver (similar to the piano fingering models in CIPI) would allow the system to predict optimal left-hand finger paths directly from clean pitch streams. This would establish a more accurate model of physical playability.
